\chapter{Descriptive Adequacy and Ontological Sufficiency}

The central argument now reaches a seam at which two questions must remain connected without being allowed to collapse into one another. The first asks which descriptive language preserves the structure of reality most faithfully. The second asks what reality is. Mathematics can improve upon ordinary language by separating distinctions that words routinely compress, and a richer mathematical formalism can resolve states that a poorer formalism treats as identical, but neither achievement is yet an ontological conclusion. A language may become more precise without thereby becoming reality itself, and a mathematical description may become more complete without demonstrating that the ontology suggested by that description is fundamental.

This distinction can be stated as a standing principle:

\begin{principle}[Descriptive adequacy is not ontological sufficiency]
A descriptive language may preserve more empirically relevant relational structure than a competing language without thereby establishing that the ontology naturally associated with that language is fundamental, complete, or uniquely true.
\end{principle}

The principle establishes the relationship between the two major levels of the inquiry. The descriptive-language analysis is logically upstream. It determines whether the mathematical vocabulary being used is sufficiently expressive to formulate the ontological problem without losing distinctions that matter. The physical--experiential analysis comes afterward and asks whether, once an adequate language has been chosen, the exterior physical description is sufficient or whether some irreducible experiential structure remains necessary. Success at the first level cannot count as evidence for a conclusion at the second. A mathematical language that resolves more distinctions than ordinary language has demonstrated greater resolving power, not a deeper ontology. A sixth-order invariant that distinguishes states quartic invariants cannot distinguish has demonstrated additional formal resolution, not an experiential dimension. A complex wave description that contains phase and directed current absent from bare geometry has demonstrated greater dynamical expressiveness, not idealism. The role of the first level is to make the second question precise enough to be tested.

The foundational persistence relation already illustrates the separation. Two reduced relational states may satisfy

\[
\bar q_t\neq\bar q_s
\]

while

\[
\Phi(\bar q_t)=\Phi(\bar q_s).
\]

The mathematical language has separated two meanings of the ordinary word ``same.'' The states are not the same, while the form is. This is a genuine gain in descriptive precision because ordinary language can move between those meanings without noticing that it has done so. The result does not establish that relational form is metaphysically more fundamental than state. It establishes that a theory of persistence that identifies sameness of state with sameness of form is unnecessarily coarse. The formalism earns its place because it allows a distinction the phenomenon requires us to make.

The same discipline governs the comparison between Platonic form and wave dynamics. Platonic symmetry supplies a language for invariant relational structure. Its group actions, projectors, characters, fixed spaces, and orbit geometry answer questions about what remains unchanged under a specified family of transformations. Wave dynamics contains many of those structural relations while adding phase, current, interference, orientation, and temporal evolution. The graph-Laplacian result makes the difference exact: one geometry can support multiple inequivalent dynamical laws, so invariant geometry cannot in general reconstruct the process that realizes it. The conclusion is that a purely geometric language is insufficient whenever the empirical phenomenon depends on phase, flow, or dynamics. It is not that wave mathematics is ontologically more real than geometry, nor that one is secretly materialistic and the other idealistic. The language-level result determines what mathematics is required to ask the physical question correctly; ontology remains to be decided by a different evidential standard.

The $\ell=5$ calculation provides the clearest example of why this separation must remain explicit. Quartic invariant structure reduces the normalized energy to a minimum condition $B_2(A)=0$. At the explicit regular state $A_\star$, the local minimum manifold has five tangent dimensions. Three are generated by rotation and therefore represent the same intrinsic form under the chosen identity relation, while two remain as genuine physical shape directions. Quartic structure cannot distinguish those two directions energetically. Degree six introduces four invariant directions not inherited from quadratic multiples of quartic invariants, and the explicit sextic $J_{4,8}$ possesses positive curvature along both shape-flat directions. The mathematical conclusion is powerful precisely because it is limited: sixth-order invariant structure resolves distinctions quartic invariant structure cannot.

The ontology-level translation is weaker by design. The result establishes that formal descriptions possess levels of resolving power and that indistinguishability under a restricted invariant algebra does not entail identity under every richer invariant algebra. It does not establish that an analogous unresolved dimension exists between physics and experience, and it certainly does not establish that phenomenal consciousness occupies any such dimension. The two quartic shape-flat directions are physical degeneracies of a truncated invariant energy. Their philosophical importance lies in demonstrating how easily a descriptive limitation can masquerade as identity when the available formal language lacks the relations required to distinguish the states.

This separation prevents the language argument from becoming deflationary. If the materialism--idealism dispute were converted into a contest over which vocabulary is more efficient, the ontological question would have been dissolved rather than answered. The physicalist claim is not merely that physical vocabulary is convenient; it is that a complete physical state is sufficient. The idealist claim is not merely that experiential language captures meanings conventional physics overlooks; it assigns explanatory or ontological priority to experience, mind, or intelligible form. A relational dual-aspect proposal likewise makes a substantive claim only if the exterior physical and interior experiential structures cannot be reduced to one another without empirical loss. These are claims about sufficiency and reality, not preferences among vocabularies.

The descriptive-language analysis therefore has a strictly limited but indispensable function: it determines whether the mathematical apparatus is capable of expressing the distinctions required to test the ontological claims. Mathematics may show that ordinary language is too coarse. Higher-order mathematics may show that an earlier mathematical description is itself too coarse. None of these results is permitted to answer the ontology question by substitution. The rule at the seam is therefore simple: descriptive adequacy determines whether we possess a language precise enough to ask the physical--experiential question; the evidential hierarchy is the only mechanism permitted to answer it.

The rule must also be applied reflexively to the present choice of relational states and invariant theory. The framework does not claim that $\mathscr R(t)$ is the uniquely correct mathematical primitive merely because relational mathematics has proven fruitful. A relational state is chosen because it is general enough to represent transformations, constraints, symmetries, dynamics, and invariant forms without assuming an object-first ontology. That is a methodological advantage, not a metaphysical proof. Category-theoretic process structures may emphasize composition rather than state identity. Causal process theories may privilege intervention and history. Measure-theoretic approaches may treat distributions rather than individual points as fundamental to prediction. Operator-algebraic approaches may encode observables without committing to a classical state ontology. The present framework earns priority only to the extent that it captures the relevant structure economically and remains testable against alternatives.

Persistence itself illustrates the need for this reflexivity. The present account represents identity largely through invariants and equivalence classes, but causal-historical continuity provides another legitimate grammar. Two states may count as one continuing system not because every defining instantaneous invariant is preserved but because they belong to one causally connected self-maintaining trajectory. Biological organisms make this especially clear. Their matter changes, their morphology changes, and many measurable quantities drift, yet causal organization, repair, memory, and regenerative continuity can preserve an identity that no single low-order snapshot invariant captures. A complete theory should therefore allow form invariance and causal continuity to complement rather than exclude one another.

The same self-skepticism applies to quotienting. Declaring two states equivalent under a group action already asserts that the corresponding difference does not matter to the scientific question. If $\Phi$ is too coarse, radically different systems can be classified as one form. If it is too fine, genuine persistence disappears beneath microscopic variation. Every important quotient and invariant must therefore be tied to an explicit identity criterion. Mathematics cannot claim superiority over ordinary language by relocating ambiguity into an unstated choice of state space, metric, group action, or form map. Formalization improves the problem when those choices themselves become inspectable and contestable.

The result is stronger than the claim that mathematics simply describes reality better than words. Mathematics becomes superior for a domain when it preserves distinctions the evidence requires, supports predictions or interventions that a coarser vocabulary cannot formulate, and exposes its assumptions sufficiently clearly for alternative representations to be compared. Ordinary language remains indispensable because empirical interpretation, phenomenological meaning, and ontological claims cannot be obtained from symbols without conceptual translation. The task is not to replace one language with another but to determine the translation relations among them and prevent success at one descriptive level from being mistaken for the truth of an ontology.

The physical--experiential question begins only after this methodological boundary has been established. Let $P_{\mathrm{phys}}$ represent an increasingly comprehensive physical state description and let $Z_\varphi$ represent a latent structure inferred from converging phenomenological evidence. The question is whether $Z_\varphi$ continues to provide predictive or interventional information after the physical description is systematically enriched. Once the relevant variables, null models, enrichment sequences, and model penalties have been preregistered, their outcomes are no longer subject to being reinterpreted away as merely another vocabulary whenever the result becomes inconvenient. The purpose of the language-level framework is precisely to construct empirical anchor points at which rival ontological programs can lose.

The same discipline governs the antiunitary bridge developed conditionally in Appendix C. The exact conjugation theorem shows that a complex wave can preserve density and symmetry while reversing current. That is a mathematical example of form surviving a transformation that reverses dynamical orientation. It becomes relevant to consciousness only if independently established phenomenological relations require the same complex projective structure. Until then, antiunitarity is an admissible mathematical hypothesis, not evidence about experience.

The governing principles can now be placed beside one another. Indistinguishability under a restricted descriptive algebra does not imply identity under every richer descriptive algebra. Descriptive-language adequacy does not imply ontological sufficiency. Mathematical admissibility does not imply phenomenal occupancy. These are not three unrelated warnings. They are the same structural discipline appearing at progressively deeper levels. A coarse invariant algebra may collapse distinctions a richer algebra resolves. An elegant mathematical language may preserve more structure without becoming an ontology. An admissible experiential representation may possess every formal property demanded of consciousness without thereby being occupied by experience. At every level, the capacity to draw a distinction must be separated from evidence that reality instantiates the thing described by that distinction.

The hierarchy is therefore fixed. First we ask whether the descriptive language is adequate to formulate the distinction. Then we ask whether the proposed mathematical structure is admissible. Then we ask whether the data occupy that structure in a way that survives prediction, intervention, cross-context transfer, and comparison with simpler alternatives. Only then are we entitled to ask what ontology best explains the result. The language of relational invariants, dynamical flow, and descriptive resolution tells us how to formulate the physical--experiential question without collapsing important distinctions. It is not allowed to answer that question in place of the evidence.
